\documentclass[11pt]{article}

% ---------- packages ----------
\usepackage[margin=1in]{geometry}
\usepackage{amsmath,amssymb,amsfonts}
\usepackage{booktabs}
\usepackage{multirow}
\usepackage{graphicx}
\usepackage{xcolor}
\usepackage{hyperref}
\usepackage{url}
\usepackage{listings}
\usepackage{caption}
\usepackage{subcaption}
\usepackage{natbib}
\usepackage{microtype}
\usepackage{enumitem}

% placeholder color
\definecolor{plhcolor}{rgb}{0.7,0.1,0.1}
\newcommand{\PH}[1][]{\textcolor{plhcolor}{\textbf{[TODO\ifx&#1&\else: #1\fi]}}}

% ---------- metadata ----------
\title{\textbf{MiniMind on H100 and MI350: A Full-Pipeline Survey of LLM\\
Training Algorithms and Inference Stacks Across Hardware Vendors}}

\author{Stephen Shao\thanks{AMD, \texttt{yu.shao@amd.com}}
\and \PH{co-authors}}

\date{\today}

\begin{document}
\maketitle

\begin{abstract}
We present a systematic survey of large language model training and inference
across two GPU platforms — NVIDIA H100 (Hopper, CUDA~12, NCCL) and AMD Instinct
MI350 (CDNA4, ROCm~7.1, RCCL) — using MiniMind, a fully from-scratch,
hand-implemented LLM pipeline in pure PyTorch, as the experimental vehicle.
MiniMind covers the full algorithm landscape from BPE tokenization and causal LM
pretraining through supervised fine-tuning (SFT), parameter-efficient LoRA
adaptation, knowledge distillation, direct preference optimization (DPO),
proximal policy optimization (PPO), group-relative policy optimization (GRPO),
and agent reinforcement learning with tool-call rollouts.  Its vendor-portable
device abstraction (\texttt{DeviceCtx}) runs the identical trainer scripts
unmodified on both hardware platforms, making the comparison about
\emph{hardware and runtime-stack differences} rather than porting effort.
We run every training stage on 1-, 4-, and 8-GPU configurations, in both dense
($\approx$64\,M parameters) and mixture-of-experts ($\approx$198\,M-A64M)
variants, and benchmark inference throughput and latency across HF-native
\texttt{generate()}, vLLM, and SGLang using Qwen3-compatible weights produced
by the in-repo conversion script.  We report wall-clock throughput,
model-FLOP utilization (MFU), peak HBM consumption, and convergence parity for
every stage.  Our results establish a precise, reproducible baseline for
vendor-neutral LLM development and highlight where stack divergence remains.
\end{abstract}

\tableofcontents
\newpage

% ---------- sections ----------
\section{Introduction}
\label{sec:intro}

Large language models (LLMs) are now trained and deployed across a rapidly
diversifying hardware landscape.  While NVIDIA GPUs and the CUDA software stack
have historically dominated LLM development, AMD Instinct accelerators running
ROCm are increasingly available in data centers and cloud offerings.  Yet most
published training results are implicitly CUDA-centric, leaving practitioners
uncertain about whether their algorithms, training frameworks, and inference
engines work correctly and efficiently on non-NVIDIA hardware.

Answering this question rigorously requires a \emph{controlled} experimental
vehicle: a codebase that implements every algorithm from scratch in portable
PyTorch, with no hidden vendor-specific dependencies, so that any observed
performance difference is attributable to the hardware and runtime stack rather
than to library porting gaps.

\paragraph{The MiniMind vehicle.}
MiniMind~\citep{minimind2024} is an end-to-end LLM pipeline implemented in pure
PyTorch without relying on \texttt{trl}, \texttt{peft}, or
\texttt{transformers.Trainer}.  It covers the full training lifecycle in a
single, coherent codebase: BPE tokenization, causal LM pretraining, supervised
fine-tuning (SFT), LoRA adaptation, knowledge distillation, direct preference
optimization (DPO), proximal policy optimization (PPO), group-relative policy
optimization (GRPO), and agent reinforcement learning with tool-call rollouts.
Its \texttt{DeviceCtx} abstraction~(\S\ref{sec:methodology}) routes all
device-specific calls — mixed-precision autocast, gradient scaling, distributed
rank binding — through a single layer that works unchanged on both CUDA and ROCm,
making it possible to run the \emph{identical} trainer scripts on both platforms.

\paragraph{Contributions.}
This paper makes the following contributions:
\begin{enumerate}[leftmargin=*]
  \item \textbf{Full-pipeline H100 vs MI350 comparison.}  We run every
    MiniMind training stage on single-node 1-, 4-, and 8-GPU configurations of
    both NVIDIA H100 (Hopper, CUDA~12, NCCL) and AMD Instinct MI350 (CDNA4,
    ROCm~7.1, RCCL), reporting throughput, model-FLOP utilization (MFU), peak
    HBM consumption, and convergence parity for each.

  \item \textbf{Dense and MoE coverage.}  We benchmark both the dense
    ($\approx$64\,M-parameter) and mixture-of-experts
    ($\approx$198\,M-A64\,M) model variants, explicitly exercising
    MoE-specific token-dispatch and router load-balancing kernels on each vendor.

  \item \textbf{Inference-stack benchmarks.}  Using the in-repo Qwen3
    conversion path, we benchmark HF-native \texttt{generate()}, vLLM, and
    SGLang on both platforms across standard workloads: latency, throughput,
    chain-of-thought reasoning (\texttt{<think>} decoding), and tool-call
    evaluation.

  \item \textbf{Cross-cutting analysis.}  We characterize convergence parity,
    \texttt{torch.compile} codegen quality, dtype maturity (bf16 / fp16 / FP8
    outlook), MoE routing behavior, and intra-node collective communication
    bandwidth (NCCL vs.\ RCCL) as cross-cutting dimensions.

  \item \textbf{A reproducible benchmark artifact.}  We release all
    experiment scripts, metric-collection utilities, and exact software version
    pins so that the community can re-run and extend these results as hardware
    and software stacks evolve.
\end{enumerate}

\paragraph{Paper organization.}
Section~\ref{sec:background} describes the MiniMind architecture, the training
algorithms covered, and the hardware/software stacks under test.
Section~\ref{sec:methodology} details our experimental methodology, including
the vendor-portable code path, reproducibility recipe, and MFU computation.
Sections~\ref{sec:training}--\ref{sec:inference} present training and inference
results respectively.  Section~\ref{sec:crosscutting} provides cross-cutting
analyses.  Sections~\ref{sec:discussion}--\ref{sec:related} discuss implications,
limitations, and related work, and Section~\ref{sec:appendix} contains the full
reproducibility appendix.

\section{Background}
\label{sec:background}

\subsection{The MiniMind Architecture}
\label{sec:background:arch}

MiniMind implements a Llama-family~\citep{touvron2023llama} transformer decoder
with the following design choices, matching the Qwen3~\citep{qwen3} model family:

\begin{itemize}[leftmargin=*]
  \item \textbf{RMSNorm} pre-normalization with $\epsilon = 10^{-6}$.
  \item \textbf{Grouped-query attention (GQA)}~\citep{ainslie2023gqa}: 8 query
    heads, 4 key-value heads, per-head RMSNorm on $Q$ and $K$.
  \item \textbf{Rotary position embeddings (RoPE)}~\citep{su2024roformer} with
    $\theta = 10^6$.  Inference-time YaRN~\citep{peng2023yarn} scaling (4$\times$
    extension) is available via \texttt{--inference\_rope\_scaling}.
  \item \textbf{SwiGLU} feed-forward network; intermediate size derived as
    $\lceil h \cdot \pi / 64 \rceil \times 64$ where $h$ is \texttt{hidden\_size}.
  \item \textbf{Flash attention}: \texttt{F.scaled\_dot\_product\_attention}
    (PyTorch SDPA), which dispatches to cuDNN-FlashAttention on CUDA and to
    hipBLASLt-fused kernels on ROCm.  Falls back to manual attention when
    \texttt{seq\_len=1} (decode step) or when a padding mask is present.
  \item \textbf{Tied embeddings} (\texttt{tie\_word\_embeddings=True});
    vocabulary size 6400, trained from a custom BPE tokenizer.
\end{itemize}

\paragraph{Dense configuration.}
Default: \texttt{hidden\_size=768}, \texttt{num\_hidden\_layers=8},
\texttt{num\_attention\_heads=8}, \texttt{num\_key\_value\_heads=4}.
Total parameters: $\approx$64\,M.

\paragraph{MoE configuration.}
A sparse mixture-of-experts~\citep{shazeer2017outrageously} variant replaces
the dense FFN with a gated MoE layer: 4 routed experts, 1 active per token, plus
a shared expert, governed by
\texttt{router\_aux\_loss\_coef}$= 5 \times 10^{-4}$.  Total parameters:
$\approx$198\,M; active per forward: $\approx$64\,M.  The loss function always
computes \texttt{loss = res.loss + res.aux\_loss} so the router
load-balancing term is included during both dense and MoE training
(it is identically zero for dense models).

\paragraph{Tokenizer.}
A BPE tokenizer with vocabulary size 6400, trained on the project's text corpus.
Special tokens include \texttt{<think>}, \texttt{</think>},
\texttt{<tool\_call>}, \texttt{<tool\_response>}, enabling chain-of-thought and
tool-use workflows directly in the chat template.

\subsection{Algorithm Landscape}
\label{sec:background:algos}

MiniMind covers the following training stages, each implemented as a self-contained
script in \texttt{trainer/}:

\paragraph{Pretraining.}
Standard causal language modeling loss over next-token prediction.  The trainer
(\texttt{train\_pretrain.py}) supports DDP via \texttt{torchrun}, bf16/fp16 mixed
precision via \texttt{torch.amp}, cosine learning-rate schedule, gradient
clipping, and optional \texttt{torch.compile}.

\paragraph{Supervised Fine-Tuning (SFT).}
\texttt{train\_full\_sft.py} uses \texttt{tokenizer.apply\_chat\_template} to
format instruction-response pairs and masks all non-assistant tokens
(labels $= -100$), so the loss is computed only on assistant completions.

\paragraph{LoRA.}
\texttt{train\_lora.py} injects low-rank adapters via a custom
\texttt{apply\_lora} function in \texttt{model/model\_lora.py} (no \texttt{peft}
dependency).  Only LoRA parameters are trained; all base-model weights are
frozen.  The adapter is saved as a separate \texttt{.pth} keyed by layer name.
\texttt{torch.compile} is automatically disabled for LoRA (monkey-patch forward
incompatibility is caught and logged).

\paragraph{Knowledge Distillation.}
\texttt{train\_distillation.py} minimizes a weighted sum of cross-entropy (CE)
loss against ground-truth labels and Kullback-Leibler (KL) divergence against
a frozen teacher's logits:
\[
  \mathcal{L} = \alpha \mathcal{L}_{\text{CE}} + (1 - \alpha)
    \cdot T^2 \cdot D_{\text{KL}}(p_{\text{teacher}} \,\|\, p_{\text{student}})
\]
where $T$ is the temperature and $\alpha \in [0,1]$ controls the CE weight.
The default configuration distills from a MoE teacher into a dense student.

\paragraph{Direct Preference Optimization (DPO).}
\texttt{train\_dpo.py} implements DPO~\citep{rafailov2023direct} with a frozen
reference model.  The loss is:
\[
  \mathcal{L}_{\text{DPO}} = -\mathbb{E}_{(x, y_w, y_l)} \left[
    \log \sigma \!\left( \beta \left(
      \log \frac{\pi_\theta(y_w|x)}{\pi_{\text{ref}}(y_w|x)} -
      \log \frac{\pi_\theta(y_l|x)}{\pi_{\text{ref}}(y_l|x)}
    \right) \right)
  \right]
\]
with $\beta = 0.15$ (default) and a very small learning rate ($4 \times 10^{-8}$)
to prevent catastrophic forgetting.

\paragraph{Proximal Policy Optimization (PPO).}
\texttt{train\_ppo.py} maintains four models simultaneously: an actor, a critic
(value model), a frozen reference, and a separately loaded reward model
(InternLM2-1.8B-reward by default).  Rollouts are collected by
\texttt{rollout\_engine.py} (\texttt{--rollout\_engine torch} or
\texttt{--rollout\_engine sglang}).  The actor is updated using clipped
PPO~\citep{schulman2017proximal} with generalized advantage estimation (GAE),
early-stopped when per-step KL exceeds a threshold.

\paragraph{Group-Relative Policy Optimization (GRPO).}
\texttt{train\_grpo.py} implements GRPO~\citep{shao2024deepseekmath}: for each
prompt, $G$ responses are sampled; rewards are computed by a reward model plus
heuristic rule terms (response length, \texttt{<think>} format, repetition
penalty); advantages are normalized within the group.  No critic model is needed.
The thinking ratio ($\texttt{--thinking\_ratio}$) controls how often the
\texttt{open\_thinking} chat-template flag is activated.

\paragraph{Agent RL.}
\texttt{train\_agent.py} extends GRPO to multi-turn agent trajectories involving
tool calls; the reward function evaluates both the quality of tool invocations and
the final answer.

\subsection{Hardware and Runtime Stacks}
\label{sec:background:hardware}

\begin{table}[h]
\centering
\caption{Hardware and software stacks under test.}
\label{tab:hw}
\small
\begin{tabular}{lll}
\toprule
\textbf{Axis} & \textbf{NVIDIA} & \textbf{AMD} \\
\midrule
GPU & H100 SXM5 (80\,GB HBM3) & MI350 (HBM3e) \\
\# GPUs / node & 8 & 8 \\
Intra-node interconnect & NVLink 4 + NVSwitch & Infinity Fabric \\
PyTorch wheel & \texttt{2.10.0+cu12x} & \texttt{2.10.0+rocm7.1} \\
Distributed backend & NCCL & RCCL (aliased as \texttt{"nccl"}) \\
Compiler stack & CUDA / cuDNN / cuBLAS / Triton & ROCm / MIOpen / hipBLASLt / Triton \\
Low-precision & bf16, fp16; Transformer Engine FP8 & bf16, fp16; ROCm FP8 \\
Inference engines & HF native, vLLM (CUDA), SGLang (CUDA) & HF native, vLLM-ROCm, SGLang-ROCm \\
\bottomrule
\end{tabular}
\end{table}

\paragraph{PyTorch portability.}
PyTorch-ROCm aliases the \texttt{torch.cuda.*} API surface to HIP so that code
referencing \texttt{device\_type="cuda"}, \texttt{torch.cuda.is\_available()},
or \texttt{torch.cuda.get\_device\_properties()} works unchanged on AMD GPUs.
MiniMind's \texttt{device\_utils.py} detects the ROCm build by inspecting
\texttt{torch.version.hip} and adjusts logging accordingly; the rest of the
codebase is unaware of the vendor.

\paragraph{SDPA dispatch.}
\texttt{F.scaled\_dot\_product\_attention} selects its backend at runtime.  On
CUDA, it may dispatch to cuDNN FlashAttention, FlashAttention-2, or
FlashAttention-3 depending on the cuDNN version and \texttt{sdp\_kernel} context.
On ROCm, it dispatches to hipBLASLt-fused attention or Composable Kernel (CK)
flash attention.  We report which backend was selected for each datapoint via
\texttt{torch.backends.cuda.flash\_sdp\_enabled()} and its ROCm equivalent.

\paragraph{Collective communication.}
Under DDP, \texttt{dist.init\_process\_group("nccl")} maps to NCCL on NVIDIA and
RCCL on AMD.  Both expose the same \texttt{torch.distributed} API.  Intra-node
bandwidth (all-reduce of the $\approx$64\,M-parameter gradient buffer at bf16)
is reported in~\S\ref{sec:crosscutting:nccl}.

\section{Methodology}
\label{sec:methodology}

\subsection{Vendor-Portable Code Path}
\label{sec:methodology:portability}

All trainer scripts use \texttt{trainer/device\_utils.py} as the single point of
contact for hardware-specific behavior.  The module exposes two layers:

\textbf{Free functions} answer environment-level questions:
\texttt{is\_rocm()}, \texttt{is\_cuda\_native()}, \texttt{vendor\_name()},
\texttt{dist\_backend()}, \texttt{default\_device\_str()},
\texttt{set\_current\_device()}, and \texttt{print\_device\_info()}.
\texttt{dist\_backend()} returns \texttt{"nccl"} on both CUDA and ROCm because
PyTorch-ROCm exposes RCCL under that string alias.
\texttt{print\_device\_info()} logs a one-line banner at startup
(Figure~\ref{fig:banner}) that is captured for every experiment run.

\textbf{\texttt{DeviceCtx}} (frozen dataclass) encapsulates per-trainer state.
Each trainer builds it once:
\begin{lstlisting}[language=Python,basicstyle=\small\ttfamily]
dev = DeviceCtx.from_arg(args.device)
if dist.is_initialized():
    dev = dev.for_rank(local_rank)
autocast_ctx = dev.autocast(args.dtype)   # nullcontext on CPU
scaler       = dev.grad_scaler(enabled=(args.dtype == 'float16'))
\end{lstlisting}
\texttt{dev.autocast()} dispatches to
\texttt{torch.amp.autocast(device\_type="cuda", ...)} for both NVIDIA and AMD
(ROCm aliases \texttt{cuda} in its PyTorch build); it returns
\texttt{nullcontext()} on CPU so no special-casing is needed.

\textbf{Consequence for this study.}  The only files that differ between a CUDA
run and a ROCm run are the PyTorch wheel and the \texttt{HSA\_OVERRIDE\_GFX\_VERSION}
environment variable (required for gfx1100-family targets; not needed for MI350
gfx940).  Trainer logic, model code, checkpoint format, and dataset loading are
bit-for-bit identical.

\subsection{Reproducibility Recipe}
\label{sec:methodology:repro}

\begin{itemize}[leftmargin=*]
  \item \textbf{Software pins.}  All experiments use:
    \texttt{torch==2.10.0+cu12x} (NVIDIA) and
    \texttt{torch==2.10.0+rocm7.1} (AMD);
    \texttt{transformers==4.5x} (pinned; see Appendix~\ref{sec:appendix:install});
    \texttt{vllm} and \texttt{sglang} at the commit SHAs listed in
    Table~\ref{tab:versions}.
  \item \textbf{Random seeds.}  Every trainer calls
    \texttt{setup\_seed(42 + rank)} at startup, which seeds Python
    \texttt{random}, NumPy, PyTorch, and \texttt{torch.cuda} (CUDA and ROCm
    use the same API surface).
  \item \textbf{Data.}  Datasets are downloaded at fixed versions; SHA-256
    hashes are recorded in Appendix~\ref{sec:appendix:data}.
  \item \textbf{Banner capture.}  Every run is launched via a wrapper that
    captures the first \texttt{[device] ...} line from
    \texttt{print\_device\_info()} and saves it alongside metrics.  This guards
    against mis-scheduled jobs.
  \item \textbf{Checkpoint naming.}  MiniMind uses load-bearing checkpoint
    names of the form
    \texttt{\{save\_weight\}\_\{hidden\_size\}\{\_moe?\}.pth}.
    The same naming convention is enforced in all experiment scripts.
\end{itemize}

\subsection{Model-FLOP Utilization}
\label{sec:methodology:mfu}

Following~\citet{chowdhery2023palm}, we define Model-FLOP Utilization (MFU) as:
\[
  \text{MFU} = \frac{\text{achieved FLOPs/s}}{\text{peak FLOPs/s of the device}}
\]
where \emph{achieved FLOPs/s} $= C \times S / T$,
with $C$ the FLOPs per token, $S$ the number of tokens processed in interval
$T$.  For a MiniMind forward + backward pass on a dense model:
\[
  C_{\text{dense}} = 6 \times P
\]
where $P$ is the number of trainable parameters (the factor of 6 accounts for
one forward and two backward passes, each costing $\approx 2P$
multiply-accumulates per token~\citep{chowdhery2023palm}).

For the MoE variant, only the active experts contribute to the compute per
token, so:
\[
  C_{\text{MoE}} = 6 \times P_{\text{active}}
\]
where $P_{\text{active}} = P_{\text{base}} + n_{\text{active}} \times
P_{\text{expert}}$, with $n_{\text{active}}$ the number of experts activated
per token (default 1 in MiniMind).

Peak FLOPs/s: H100 SXM5 delivers 989\,TFLOPs/s (bf16 tensor core);
MI350 delivers \PH{peak TFLOPs/s from AMD spec sheet}.
MFU is computed offline from logged step times by
\texttt{paper/experiments/compute\_mfu.py}.

\subsection{Inference Metrics}
\label{sec:methodology:inf}

\begin{itemize}[leftmargin=*]
  \item \textbf{TTFT} (time-to-first-token): wall time from request arrival to
    first output token, measured at the server level for
    \texttt{serve\_openai\_api.py} and at the engine level for vLLM / SGLang.
  \item \textbf{Inter-token latency} (ITL): wall time between consecutive output
    tokens during decode; reported as p50 and p95 over a 200-request run.
  \item \textbf{Tokens/s}: aggregate decode throughput across concurrent
    requests.
  \item \textbf{Requests/s}: completed requests per second at each concurrency
    level.
  \item \textbf{HBM at steady state}: \texttt{torch.cuda.max\_memory\_allocated()}
    sampled once per second during the benchmark window.
  \item \textbf{Prefill/decode split}: derived by logging TTFT separately from
    total generation time.
\end{itemize}

\subsection{Experimental Configurations}
\label{sec:methodology:configs}

Table~\ref{tab:configs} summarizes the experimental configurations used
throughout this paper.  For training, we use the ``mini'' datasets shipped with
the MiniMind repository, which are sufficient to reproduce the upstream
single-3090 reference numbers and yield comparable relative performance across
platforms.  Full dataset hashes are in Appendix~\ref{sec:appendix:data}.

\begin{table}[h]
\centering
\caption{Summary of experimental configurations.}
\label{tab:configs}
\small
\begin{tabular}{llll}
\toprule
\textbf{Dimension} & \textbf{Values} & \textbf{Primary} & \textbf{Note} \\
\midrule
Model variant & Dense, MoE & Dense & hidden\_size=768, 8 layers \\
\# GPUs & 1, 4, 8 & 8 & Single node, DDP \\
dtype & bf16, fp16 & bf16 & fp16 uses GradScaler \\
\texttt{torch.compile} & 0, 1 & 0 & LoRA: auto-disabled \\
Rollout engine & torch, sglang & torch & PPO/GRPO only \\
Inference engine & HF, vLLM, SGLang & vLLM & Block E \\
\bottomrule
\end{tabular}
\end{table}

\section{Training-Stage Results}
\label{sec:training}

\subsection{Block A: Pretraining}
\label{sec:training:pretrain}

Table~\ref{tab:pretrain_dense} reports pretraining throughput and MFU for the
dense configuration across GPU counts and vendors.
Table~\ref{tab:pretrain_moe} gives the corresponding MoE results.

\begin{table}[h]
\centering
\caption{Dense pretrain (hidden\_size=768, 8 layers, bf16, 2 epochs, pretrain\_t2t\_mini).
MFU computed against H100 989\,TFLOPs/s (NVIDIA) and MI350 \PH{X}\,TFLOPs/s (AMD).}
\label{tab:pretrain_dense}
\small
\begin{tabular}{lllrrrrr}
\toprule
\textbf{Vendor} & \textbf{GPUs} & \textbf{dtype / compile} &
\textbf{tok/s} & \textbf{MFU (\%)} & \textbf{Peak HBM (GB)} &
\textbf{Epoch (min)} & \textbf{Final loss} \\
\midrule
H100  & 1 & bf16 / off & \PH{} & \PH{} & \PH{} & \PH{} & \PH{} \\
H100  & 4 & bf16 / off & \PH{} & \PH{} & \PH{} & \PH{} & \PH{} \\
H100  & 8 & bf16 / off & \PH{} & \PH{} & \PH{} & \PH{} & \PH{} \\
H100  & 8 & bf16 / on  & \PH{} & \PH{} & \PH{} & \PH{} & \PH{} \\
H100  & 8 & fp16 / off & \PH{} & \PH{} & \PH{} & \PH{} & \PH{} \\
\midrule
MI350 & 1 & bf16 / off & \PH{} & \PH{} & \PH{} & \PH{} & \PH{} \\
MI350 & 4 & bf16 / off & \PH{} & \PH{} & \PH{} & \PH{} & \PH{} \\
MI350 & 8 & bf16 / off & \PH{} & \PH{} & \PH{} & \PH{} & \PH{} \\
MI350 & 8 & bf16 / on  & \PH{} & \PH{} & \PH{} & \PH{} & \PH{} \\
MI350 & 8 & fp16 / off & \PH{} & \PH{} & \PH{} & \PH{} & \PH{} \\
\bottomrule
\end{tabular}
\end{table}

\begin{table}[h]
\centering
\caption{MoE pretrain (4 experts, 1 active, hidden\_size=768, bf16). MFU uses $P_{\text{active}}$.}
\label{tab:pretrain_moe}
\small
\begin{tabular}{llrrrrr}
\toprule
\textbf{Vendor} & \textbf{GPUs} &
\textbf{tok/s} & \textbf{MFU (\%)} & \textbf{Peak HBM (GB)} &
\textbf{Epoch (min)} & \textbf{Final loss + aux} \\
\midrule
H100  & 1 & \PH{} & \PH{} & \PH{} & \PH{} & \PH{} \\
H100  & 4 & \PH{} & \PH{} & \PH{} & \PH{} & \PH{} \\
H100  & 8 & \PH{} & \PH{} & \PH{} & \PH{} & \PH{} \\
\midrule
MI350 & 1 & \PH{} & \PH{} & \PH{} & \PH{} & \PH{} \\
MI350 & 4 & \PH{} & \PH{} & \PH{} & \PH{} & \PH{} \\
MI350 & 8 & \PH{} & \PH{} & \PH{} & \PH{} & \PH{} \\
\bottomrule
\end{tabular}
\end{table}

\paragraph{Multi-GPU scaling.}
Figure~\ref{fig:scaling} (placeholder) shows strong-scaling efficiency
(normalized throughput per GPU vs.\ ideal linear) for both platforms.
We expect near-linear scaling within 8 GPUs for a 64\,M-parameter model
whose activation size is small relative to intra-node interconnect bandwidth.

\paragraph{\texttt{torch.compile} effect.}
\PH{Discuss compile speedup on H100 vs MI350 once data are available.}

\subsection{Block B: Supervised Fine-Tuning, LoRA, and Distillation}
\label{sec:training:posttrain}

\begin{table}[h]
\centering
\caption{Post-training stages: throughput and memory. SFT uses seq\_len=768,
LoRA trains only adapter parameters ($\ll 1\%$ of total), Distillation
runs MoE teacher + dense student simultaneously.}
\label{tab:posttrain}
\small
\begin{tabular}{lllrrr}
\toprule
\textbf{Stage} & \textbf{Vendor} & \textbf{GPUs} &
\textbf{tok/s} & \textbf{Peak HBM (GB)} & \textbf{Final loss} \\
\midrule
SFT         & H100  & 1 & \PH{} & \PH{} & \PH{} \\
SFT         & H100  & 8 & \PH{} & \PH{} & \PH{} \\
SFT         & MI350 & 1 & \PH{} & \PH{} & \PH{} \\
SFT         & MI350 & 8 & \PH{} & \PH{} & \PH{} \\
\midrule
LoRA        & H100  & 1 & \PH{} & \PH{} & \PH{} \\
LoRA        & H100  & 8 & \PH{} & \PH{} & \PH{} \\
LoRA        & MI350 & 1 & \PH{} & \PH{} & \PH{} \\
LoRA        & MI350 & 8 & \PH{} & \PH{} & \PH{} \\
\midrule
Distill (MoE$\to$dense) & H100  & 8 & \PH{} & \PH{} & \PH{} \\
Distill (MoE$\to$dense) & MI350 & 8 & \PH{} & \PH{} & \PH{} \\
\bottomrule
\end{tabular}
\end{table}

\paragraph{LoRA efficiency.}
Because \texttt{torch.compile} is automatically disabled for LoRA (the
monkey-patched forward is incompatible; see \texttt{train\_lora.py:165}),
the throughput gap vs.\ full SFT is entirely a function of trainable-parameter
ratio, not kernel compilation.

\paragraph{Distillation with a MoE teacher.}
The distillation step exercises MoE forward passes in the teacher alongside
dense backward passes in the student on every step, stressing memory bandwidth
differently from pure-MoE or pure-dense training.
\PH{Discuss observed memory and throughput differences between vendors.}

\subsection{Block C: Preference Alignment (DPO, PPO, GRPO)}
\label{sec:training:rlhf}

\begin{table}[h]
\centering
\caption{DPO results. The frozen reference doubles activation memory;
 batch size is limited to 4 per GPU.}
\label{tab:dpo}
\small
\begin{tabular}{llrrrr}
\toprule
\textbf{Vendor} & \textbf{GPUs} &
\textbf{tok/s} & \textbf{Peak HBM (GB)} &
\textbf{DPO loss} & \textbf{Epoch (min)} \\
\midrule
H100  & 1 & \PH{} & \PH{} & \PH{} & \PH{} \\
H100  & 8 & \PH{} & \PH{} & \PH{} & \PH{} \\
MI350 & 1 & \PH{} & \PH{} & \PH{} & \PH{} \\
MI350 & 8 & \PH{} & \PH{} & \PH{} & \PH{} \\
\bottomrule
\end{tabular}
\end{table}

\begin{table}[h]
\centering
\caption{GRPO and PPO rollout throughput and training metrics.
Rollout-tok/s is measured inside the training loop via the \texttt{RolloutResult} timer.}
\label{tab:rlhf}
\small
\begin{tabular}{lllrrrrr}
\toprule
\textbf{Algo} & \textbf{Vendor} & \textbf{Engine} &
\textbf{$G$} & \textbf{rollout-tok/s} & \textbf{train-tok/s} &
\textbf{Peak HBM (GB)} & \textbf{Reward} \\
\midrule
GRPO & H100  & torch  & 4 & \PH{} & \PH{} & \PH{} & \PH{} \\
GRPO & H100  & torch  & 8 & \PH{} & \PH{} & \PH{} & \PH{} \\
GRPO & H100  & sglang & 4 & \PH{} & \PH{} & \PH{} & \PH{} \\
GRPO & MI350 & torch  & 4 & \PH{} & \PH{} & \PH{} & \PH{} \\
GRPO & MI350 & torch  & 8 & \PH{} & \PH{} & \PH{} & \PH{} \\
GRPO & MI350 & sglang & 4 & \PH{} & \PH{} & \PH{} & \PH{} \\
\midrule
PPO  & H100  & torch  & -- & \PH{} & \PH{} & \PH{} & \PH{} \\
PPO  & H100  & sglang & -- & \PH{} & \PH{} & \PH{} & \PH{} \\
PPO  & MI350 & torch  & -- & \PH{} & \PH{} & \PH{} & \PH{} \\
PPO  & MI350 & sglang & -- & \PH{} & \PH{} & \PH{} & \PH{} \\
\bottomrule
\end{tabular}
\end{table}

\subsection{Block D: Agent RL}
\label{sec:training:agent}

Table~\ref{tab:agent} reports trajectory throughput and tool-call success rate.
Because agent training involves multi-turn rollouts with real tool invocations,
the rollout step is significantly more variable than in GRPO/PPO.

\begin{table}[h]
\centering
\caption{Agent RL: trajectory throughput and tool-call success.}
\label{tab:agent}
\small
\begin{tabular}{llrrrr}
\toprule
\textbf{Vendor} & \textbf{GPUs} &
\textbf{traj/h} & \textbf{tool-call success (\%)} &
\textbf{Peak HBM (GB)} & \textbf{Reward curve} \\
\midrule
H100  & 1 & \PH{} & \PH{} & \PH{} & \PH{} \\
H100  & 8 & \PH{} & \PH{} & \PH{} & \PH{} \\
MI350 & 1 & \PH{} & \PH{} & \PH{} & \PH{} \\
MI350 & 8 & \PH{} & \PH{} & \PH{} & \PH{} \\
\bottomrule
\end{tabular}
\end{table}

\section{Inference-Engine Results}
\label{sec:inference}

All inference benchmarks use weights produced by
\texttt{scripts/convert\_model.py} (\texttt{convert\_torch2transformers}),
which exports native \texttt{.pth} checkpoints to \texttt{Qwen3ForCausalLM} /
\texttt{Qwen3MoeForCausalLM} HuggingFace format compatible with vLLM and SGLang.
Greedy-output parity between the native MiniMind model and each engine was
verified as Gate~6 (\S\ref{sec:methodology}).

\subsection{HF-Native \texttt{generate()}}
\label{sec:inference:hf}

\begin{table}[h]
\centering
\caption{HF-native single-stream decode throughput (tokens/s). Single GPU,
\texttt{model.half().eval().to(device)}, temperature=0.85, top\_p=0.95.}
\label{tab:hf_single}
\small
\begin{tabular}{llrrr}
\toprule
\textbf{Model} & \textbf{Vendor} & \textbf{128 tok} & \textbf{512 tok} & \textbf{2048 tok} \\
\midrule
Dense & H100  & \PH{} & \PH{} & \PH{} \\
Dense & MI350 & \PH{} & \PH{} & \PH{} \\
MoE   & H100  & \PH{} & \PH{} & \PH{} \\
MoE   & MI350 & \PH{} & \PH{} & \PH{} \\
\bottomrule
\end{tabular}
\end{table}

\begin{table}[h]
\centering
\caption{HF-native concurrent throughput via \texttt{serve\_openai\_api.py}
(dense model, max\_tokens=512, 200 requests).}
\label{tab:hf_concurrent}
\small
\begin{tabular}{llrrrr}
\toprule
\textbf{Vendor} & \textbf{Concurrency} &
\textbf{TTFT p50 (ms)} & \textbf{TTFT p95 (ms)} &
\textbf{tok/s} & \textbf{req/s} \\
\midrule
H100  & 1   & \PH{} & \PH{} & \PH{} & \PH{} \\
H100  & 8   & \PH{} & \PH{} & \PH{} & \PH{} \\
H100  & 32  & \PH{} & \PH{} & \PH{} & \PH{} \\
H100  & 128 & \PH{} & \PH{} & \PH{} & \PH{} \\
\midrule
MI350 & 1   & \PH{} & \PH{} & \PH{} & \PH{} \\
MI350 & 8   & \PH{} & \PH{} & \PH{} & \PH{} \\
MI350 & 32  & \PH{} & \PH{} & \PH{} & \PH{} \\
MI350 & 128 & \PH{} & \PH{} & \PH{} & \PH{} \\
\bottomrule
\end{tabular}
\end{table}

\subsection{vLLM}
\label{sec:inference:vllm}

\begin{table}[h]
\centering
\caption{vLLM throughput sweep (dense Qwen3-converted weights, 8 GPUs,
tensor\_parallel\_size=1, max\_tokens=512).}
\label{tab:vllm}
\small
\begin{tabular}{llrrrr}
\toprule
\textbf{Vendor} & \textbf{Concurrency} &
\textbf{TTFT p50 (ms)} & \textbf{ITL p50 (ms)} &
\textbf{tok/s} & \textbf{req/s} \\
\midrule
H100  & 1   & \PH{} & \PH{} & \PH{} & \PH{} \\
H100  & 8   & \PH{} & \PH{} & \PH{} & \PH{} \\
H100  & 32  & \PH{} & \PH{} & \PH{} & \PH{} \\
H100  & 128 & \PH{} & \PH{} & \PH{} & \PH{} \\
\midrule
MI350 & 1   & \PH{} & \PH{} & \PH{} & \PH{} \\
MI350 & 8   & \PH{} & \PH{} & \PH{} & \PH{} \\
MI350 & 32  & \PH{} & \PH{} & \PH{} & \PH{} \\
MI350 & 128 & \PH{} & \PH{} & \PH{} & \PH{} \\
\bottomrule
\end{tabular}
\end{table}

\subsection{SGLang}
\label{sec:inference:sglang}

\begin{table}[h]
\centering
\caption{SGLang throughput sweep (dense Qwen3-converted weights, max\_tokens=512).}
\label{tab:sglang}
\small
\begin{tabular}{llrrrr}
\toprule
\textbf{Vendor} & \textbf{Concurrency} &
\textbf{TTFT p50 (ms)} & \textbf{ITL p50 (ms)} &
\textbf{tok/s} & \textbf{req/s} \\
\midrule
H100  & 1   & \PH{} & \PH{} & \PH{} & \PH{} \\
H100  & 32  & \PH{} & \PH{} & \PH{} & \PH{} \\
H100  & 128 & \PH{} & \PH{} & \PH{} & \PH{} \\
\midrule
MI350 & 1   & \PH{} & \PH{} & \PH{} & \PH{} \\
MI350 & 32  & \PH{} & \PH{} & \PH{} & \PH{} \\
MI350 & 128 & \PH{} & \PH{} & \PH{} & \PH{} \\
\bottomrule
\end{tabular}
\end{table}

\subsection{Reasoning and Tool-Call Workloads}
\label{sec:inference:reasoning}

\begin{table}[h]
\centering
\caption{Chain-of-thought (\texttt{open\_thinking=True}) decode throughput
vs.\ standard decode (concurrency=1, max\_tokens=2048, HF-native).}
\label{tab:thinking}
\small
\begin{tabular}{llrr}
\toprule
\textbf{Vendor} & \textbf{Mode} & \textbf{tok/s (decode only)} & \textbf{Think tokens (\%)} \\
\midrule
H100  & Standard  & \PH{} & --- \\
H100  & Thinking  & \PH{} & \PH{} \\
MI350 & Standard  & \PH{} & --- \\
MI350 & Thinking  & \PH{} & \PH{} \\
\bottomrule
\end{tabular}
\end{table}

\begin{table}[h]
\centering
\caption{Tool-call evaluation (\texttt{scripts/eval\_toolcall.py}): success rate and latency.}
\label{tab:toolcall}
\small
\begin{tabular}{lrr}
\toprule
\textbf{Vendor} & \textbf{Tool-call success (\%)} & \textbf{Latency p50 (ms)} \\
\midrule
H100  & \PH{} & \PH{} \\
MI350 & \PH{} & \PH{} \\
\bottomrule
\end{tabular}
\end{table}

\section{Cross-Cutting Analyses}
\label{sec:crosscutting}

\subsection{Convergence Parity}
\label{sec:crosscutting:convergence}

A prerequisite for attributing throughput differences to hardware rather than
numerical divergence is that both platforms converge to the same loss trajectory
under identical seeds.  We verify this via Gate~2 of the verification harness
(\S\ref{sec:methodology}): a single-GPU dense pretrain is run on each vendor
with \texttt{setup\_seed(42)} and bf16 AMP; the cross-entropy loss at step~100
must agree within $\pm 2\%$ relative.

\begin{table}[h]
\centering
\caption{Convergence parity (dense pretrain, bf16, 1 GPU, seed=42).
$\delta_{\text{rel}} = |L_{\text{H100}} - L_{\text{MI350}}| / \max(L_{\text{H100}}, L_{\text{MI350}})$.
Tolerance: $\leq 2\%$.}
\label{tab:convergence}
\small
\begin{tabular}{lrrrr}
\toprule
\textbf{Step} & \textbf{H100 loss} & \textbf{MI350 loss} & $\delta_{\text{rel}}$ (\%) & \textbf{Pass?} \\
\midrule
100  & \PH{} & \PH{} & \PH{} & \PH{} \\
500  & \PH{} & \PH{} & \PH{} & \PH{} \\
1000 & \PH{} & \PH{} & \PH{} & \PH{} \\
\midrule
Final (epoch 2) & \PH{} & \PH{} & \PH{} & \PH{} \\
\bottomrule
\end{tabular}
\end{table}

PyTorch-ROCm aliases the \texttt{torch.cuda.*} namespace and uses the same
cuBLAS-equivalent kernel dispatch path for GEMM workloads via hipBLASLt, so
convergence parity is expected in bf16.  Any deviation beyond 2\% would
indicate a fused-kernel discrepancy (e.g.\ in the SwiGLU or RMSNorm fused
backward) and would require per-operator investigation before the
throughput numbers can be meaningfully compared.

For the MoE configuration, we additionally verify that the router's soft-max
probabilities follow the same distribution at initialization (both vendors use
PyTorch's default Kaiming-uniform initialization, which must be seeded
identically).  The \texttt{aux\_loss} trajectory is a proxy for router load
balance and is reported alongside cross-entropy loss in Table~\ref{tab:convergence_moe}.

\begin{table}[h]
\centering
\caption{MoE convergence parity (bf16, 1 GPU, seed=42).  Both CE loss and
auxiliary load-balance loss are compared.}
\label{tab:convergence_moe}
\small
\begin{tabular}{lrrrr}
\toprule
\textbf{Step} & \textbf{H100 CE / aux} & \textbf{MI350 CE / aux} & $\delta_{\text{rel}}^{\text{CE}}$ (\%) & \textbf{Pass?} \\
\midrule
100  & \PH{} / \PH{} & \PH{} / \PH{} & \PH{} & \PH{} \\
500  & \PH{} / \PH{} & \PH{} / \PH{} & \PH{} & \PH{} \\
1000 & \PH{} / \PH{} & \PH{} / \PH{} & \PH{} & \PH{} \\
\bottomrule
\end{tabular}
\end{table}

\subsection{dtype Maturity: bf16, fp16, and the FP8 Outlook}
\label{sec:crosscutting:dtype}

\paragraph{bf16 (primary).}
Both H100 and MI350 natively support bf16 tensor cores.  MiniMind trainers
default to bf16 via \texttt{torch.amp.autocast(device\_type=device, dtype=torch.bfloat16)}.
The GradScaler is instantiated but effectively a no-op for bf16 (loss scaling
is not needed due to the wider dynamic range).  We observe numerically stable
training on both platforms with no manual loss-scale intervention.

\paragraph{fp16 (secondary).}
MiniMind supports fp16 via the same AMP path with GradScaler enabled.  fp16
is more sensitive to gradient underflow, particularly in the early warmup
phase of pretraining.  Table~\ref{tab:pretrain_dense} includes fp16 rows so
readers can verify that neither vendor exhibits systematic instability beyond
the known GradScaler overhead.

\paragraph{FP8 (future).}
Neither the current training scripts nor \texttt{torch.amp} expose a stable
FP8 path in PyTorch~2.10.  NVIDIA exposes FP8 via Transformer Engine (TE);
AMD exposes a parallel path via ROCm~7.x \texttt{hipBLASLt} FP8 GEMM
routines.  Because MiniMind's hand-written layers bypass TE's
\texttt{Linear} wrappers, enabling FP8 would require manual \texttt{Float8Linear}
substitution and is left as future work.  We flag in
Table~\ref{tab:dtype_matrix} which dtypes were exercised and which are
expected to become available.

\begin{table}[h]
\centering
\caption{dtype availability and training-stability status (MiniMind~2.10 baseline).}
\label{tab:dtype_matrix}
\small
\begin{tabular}{lllll}
\toprule
\textbf{dtype} & \textbf{H100 support} & \textbf{MI350 support} & \textbf{Trainer path} & \textbf{Status} \\
\midrule
bf16  & \checkmark & \checkmark & \texttt{autocast bf16 + no-op scaler} & Stable \\
fp16  & \checkmark & \checkmark & \texttt{autocast fp16 + GradScaler}   & Stable \\
fp32  & \checkmark & \checkmark & \texttt{nullcontext (default CPU)}    & Baseline only \\
FP8   & TE path    & ROCm 7.x   & Not implemented                      & Future \\
\bottomrule
\end{tabular}
\end{table}

\subsection{MoE Routing Behavior}
\label{sec:crosscutting:moe}

MiniMind's MoE FFN routes each token to $k{=}1$ of $N{=}4$ routed experts
plus a shared expert that always runs.  The router is a linear projection
followed by softmax; the top-$k$ gate values weight the expert outputs.
A load-balancing auxiliary loss $\mathcal{L}_{\text{aux}} = \alpha \sum_i f_i \cdot P_i$
(fraction of tokens $f_i$ vs.\ mean gate probability $P_i$ per expert) is
added to the primary CE loss at every step.

The key vendor-specific concern is the \emph{token-dispatch scatter/gather}
pattern: on CUDA, expert computation can be fused using custom Triton kernels
(or cuBLAS batched-GEMM with grouped indices); on ROCm, the same
\texttt{torch.index\_select} / \texttt{scatter\_add} path is used, falling
back to hipBLASLt for the GEMM portion.  For MiniMind's small expert count
($N{=}4$) and small hidden dimension (768), the dispatch overhead is
negligible relative to the GEMM time, but we track it explicitly via
step-time decomposition.

\begin{table}[h]
\centering
\caption{MoE expert utilization at epoch 1 end (pretrain, bf16, 1 GPU).
An ideal load-balanced router assigns $25\%$ of tokens to each routed expert.
Utilization $> 5\%$ deviation from $25\%$ may indicate router collapse.}
\label{tab:moe_balance}
\small
\begin{tabular}{lrrrrrr}
\toprule
\textbf{Vendor} & \textbf{Expert 0 (\%)} & \textbf{Expert 1 (\%)} &
\textbf{Expert 2 (\%)} & \textbf{Expert 3 (\%)} &
\textbf{Shared (\%)} & \textbf{aux\_loss} \\
\midrule
H100  & \PH{} & \PH{} & \PH{} & \PH{} & 100 & \PH{} \\
MI350 & \PH{} & \PH{} & \PH{} & \PH{} & 100 & \PH{} \\
\bottomrule
\end{tabular}
\end{table}

\subsection{\texttt{torch.compile} Codegen Quality}
\label{sec:crosscutting:compile}

\texttt{torch.compile} is exercised via the \texttt{--use\_compile 1} flag
available in all MiniMind trainers.  The compiler backend on CUDA (Ampere+)
is Triton/inductor, which generates efficient CUDA kernels for the attention
SDPA, layer norms, and element-wise operations.  On ROCm, the same
\texttt{torch.compile} path uses the Triton-ROCm backend (upstreamed into
PyTorch~2.3+), which emits LLVM-rooted HIP kernels via HSACO.

Known limitation: \texttt{torch.compile} is automatically disabled for LoRA
runs (\texttt{train\_lora.py:165}) because the monkey-patched
\texttt{Linear.forward} is not graph-capture-compatible.  This limitation
applies equally on both vendors.

\begin{table}[h]
\centering
\caption{\texttt{torch.compile} speedup over eager mode (pretrain, dense,
bf16, 8 GPUs, epoch 1 throughput).  Speedup is computed as
$\text{tok/s}_{\text{compile}} / \text{tok/s}_{\text{eager}} - 1$.}
\label{tab:compile}
\small
\begin{tabular}{lrrrr}
\toprule
\textbf{Vendor} & \textbf{Eager tok/s} & \textbf{Compile tok/s} & \textbf{Speedup (\%)} & \textbf{Compile time (s)} \\
\midrule
H100  & \PH{} & \PH{} & \PH{} & \PH{} \\
MI350 & \PH{} & \PH{} & \PH{} & \PH{} \\
\bottomrule
\end{tabular}
\end{table}

\PH{Discuss observed compile-time and warm-up overhead on each vendor once
data are available.  On AMD, the first-run Triton-ROCm kernel compilation
can take 5--15 minutes for a new model shape, which inflates the first-epoch
wall time; subsequent runs use the HSACO cache (\texttt{\$HOME/.triton/cache}).}

\subsection{NCCL vs.\ RCCL: DDP All-Reduce Bandwidth}
\label{sec:crosscutting:collective}

MiniMind uses PyTorch DDP, which performs one all-reduce per parameter tensor
per backward pass.  For a 64M-parameter dense model in bf16, the total
gradient buffer size is $64{\times}10^6 \times 2\,\text{B} \approx 128\,\text{MB}$
per all-reduce.  For the 198M-A64M MoE model (total parameters, bf16), the
buffer is $\approx 396\,\text{MB}$.

On the H100 side, NCCL uses NVLink~4 (900\,GB/s bidirectional per GPU pair
via NVSwitch) for intra-node all-reduce; effective ring-allreduce bandwidth
scales with the number of GPUs up to the bisection bandwidth limit.  On the
MI350 side, RCCL uses AMD Infinity Fabric (the specific bandwidth for MI350
in an 8-GPU OAM configuration is \PH{}) exposed as the \texttt{"nccl"} backend
(RCCL is API-compatible with NCCL and is aliased transparently by
\texttt{dist\_backend()} in \texttt{trainer/device\_utils.py}).

\begin{table}[h]
\centering
\caption{DDP all-reduce wall time for a 128\,MB gradient buffer (bf16, 8 GPUs,
bus-bandwidth measurement via \texttt{nccl-tests} or equivalent).}
\label{tab:allreduce}
\small
\begin{tabular}{lrrrr}
\toprule
\textbf{Vendor} & \textbf{Backend} & \textbf{Latency ($\mu$s)} &
\textbf{Bandwidth (GB/s)} & \textbf{Peak (\%)} \\
\midrule
H100  & NCCL  & \PH{} & \PH{} & \PH{} \\
MI350 & RCCL  & \PH{} & \PH{} & \PH{} \\
\bottomrule
\end{tabular}
\end{table}

Because MiniMind's model is small relative to the intra-node interconnect
bandwidth on either platform, we expect all-reduce to be a sub-dominant
cost ($< 5\%$ of step time at 8 GPUs) for both vendors.  This simplifies
the strong-scaling analysis in \S\ref{sec:training:pretrain}: near-linear
scaling is driven by compute efficiency, not communication overhead.

\section{Discussion}
\label{sec:discussion}

\subsection{The Floor of Vendor-Neutral LLM Development}
\label{sec:discussion:floor}

MiniMind demonstrates that the \emph{algorithmic} layer of LLM development
is largely vendor-neutral today.  Every training stage — from BPE tokenizer
through agent RL — runs on both H100 and MI350 without changes to the
trainer code, using only the portability layer in \texttt{trainer/device\_utils.py}
(\texttt{DeviceCtx}, \texttt{dist\_backend()}, \texttt{default\_device\_str()}).
PyTorch-ROCm's aliasing of the \texttt{torch.cuda.*} surface is sufficiently
mature that the remaining divergences are confined to:

\begin{enumerate}
\item \textbf{SDPA backend selection.} \texttt{F.scaled\_dot\_product\_attention}
  dispatches to cuDNN FlashAttention (CUDA) or to hipBLASLt-fused attention /
  Composable Kernel (ROCm) depending on the runtime.  The same PyTorch API
  surface is used in both cases; the performance and numerical properties of
  the underlying kernels differ.
\item \textbf{\texttt{torch.compile} warm-up time.}  First-run Triton-ROCm
  kernel compilation is slower than CUDA due to the additional LLVM→HSACO
  step.  Subsequent runs are cached.
\item \textbf{FP8 availability.}  NVIDIA exposes FP8 via Transformer Engine
  with a stable API; AMD's FP8 path via hipBLASLt is available but not yet
  wrapped in a \texttt{torch.nn.Linear}-compatible interface in stock PyTorch.
\item \textbf{Inference ecosystem maturity.}  vLLM and SGLang maintain
  ROCm branches, but the feature gap (e.g.\ speculative decoding, RadixAttention
  prefix caching) between the CUDA-primary and ROCm paths closes with each
  release cycle.
\end{enumerate}

\subsection{Where Divergence Remains}
\label{sec:discussion:divergence}

\paragraph{Custom kernels.}
Projects that rely on hand-written CUDA kernels (Flash-Attention~3, DeepSpeed
Triton extensions, Liger-Kernel) face a porting cost on ROCm.  MiniMind
deliberately avoids this by relying exclusively on \texttt{F.scaled\_dot\_product\_attention}
and standard PyTorch operations.  This is a deliberate design choice that
trades potential peak performance for portability.

\paragraph{MoE token dispatch.}
Expert routing involves a scatter/gather pattern that benefits from
vendor-specific fused kernels.  On CUDA, libraries such as Megablocks or
cuBLAS batched-GEMM with custom dispatch achieve near-zero overhead.
On ROCm, the same functionality is available via Composable Kernel (CK) and
hipBLASLt grouped GEMM, but the integration with PyTorch's dispatcher is
less mature.  For MiniMind's small ($N{=}4$) expert count and 768-dim hidden
size, the overhead is negligible, but larger MoE models would expose this
gap more sharply.

\paragraph{Serving ecosystem.}
vLLM and SGLang are CUDA-first; the ROCm ports exist but lag on features
such as FP8 KV-cache quantization, speculative decoding, and disaggregated
prefill-decode.  For production serving, these gaps translate to a real
efficiency penalty on AMD hardware that is separate from raw compute throughput.

\paragraph{Tooling and profiling.}
NVIDIA Nsight Systems/Compute provides fine-grained kernel-level profiling
that is tightly integrated with the CUDA runtime.  AMD's equivalent
(Omniperf, ROCprofiler-SDK, RocTracer) has reached parity for most common
workloads but has a steeper learning curve and less community-written tooling.

\subsection{Implications for the Research Community}
\label{sec:discussion:implications}

Our results provide three concrete takeaways:

\begin{enumerate}
\item \textbf{Algorithm development can and should be vendor-neutral.}
  The overhead of maintaining a portable device layer
  (\texttt{device\_utils.py}: 120 lines) is small relative to the algorithmic
  work in each trainer.  New algorithm proposals should publish on
  vendor-neutral code paths to avoid creating AMD-inaccessible baselines.

\item \textbf{The inference gap is the primary remaining concern.}
  Training throughput parity (where it exists) does not imply serving parity.
  The vLLM and SGLang ROCm feature gaps mean that AMD hardware may train
  efficiently but serve at reduced throughput until the inference ecosystem
  catches up.

\item \textbf{Small models are ideal for cross-vendor evaluation.}
  MiniMind's 64M parameters fit easily in a single-GPU context, which
  eliminates FSDP/TP parallelism as a confound.  Researchers validating
  vendor portability of new algorithms should prototype on small models first
  before investing in multi-node, mixed-precision, or expert-parallel stacks.
\end{enumerate}

\section{Limitations}
\label{sec:limitations}

\paragraph{Single-node scope.}
All experiments run on a single 8-GPU node.  Multi-node training introduces
InfiniBand or RoCE interconnects, where NCCL and RCCL may diverge significantly
in all-reduce latency and in the choice of collective algorithms (tree vs.\
ring vs.\ NVLS).  Our conclusions about communication overhead apply only to
the intra-node (NVLink / Infinity Fabric) regime.

\paragraph{Small model scale.}
MiniMind is a 64M-parameter dense / 198M-A64M MoE model.  At this scale,
the model fits comfortably within a single GPU's HBM, and DDP all-reduce
represents a small fraction of step time.  Larger models (7B+) would expose
memory-bandwidth bottlenecks, activation-checkpoint tradeoffs, and
collective-communication overheads that are not visible here.  MFU numbers
at small scale tend to be lower than at large scale (arithmetic intensity is
suboptimal), so our absolute MFU values should not be extrapolated directly
to production-scale training.

\paragraph{One model family.}
MiniMind's architecture is a Llama-family model (RMSNorm, RoPE, GQA, SwiGLU).
Different architectural choices — state-space models (Mamba), linear attention
variants, or convolutional hybrids — may exhibit different vendor-portability
profiles, particularly for operations not covered by \texttt{F.scaled\_dot\_product\_attention}.

\paragraph{Specific software pins.}
All measurements use PyTorch~2.10, CUDA~12.x, and ROCm~7.1 (see
Appendix~\ref{sec:appendix} for exact pins).  Both stacks evolve rapidly;
the ROCm performance gap on SDPA and \texttt{torch.compile} warm-up is
expected to narrow with each PyTorch release.  Results should be revalidated
when upgrading software.

\paragraph{Dataset scale.}
We use the ``mini'' dataset variants shipped with the MiniMind repository
(\texttt{pretrain\_t2t\_mini}, \texttt{sft\_t2t\_mini}).  These are
designed for rapid iteration, not state-of-the-art model quality.  Loss
and perplexity values reported here reflect training on small corpora and
should not be compared to results trained on web-scale data.

\paragraph{Inference engine version skew.}
The ROCm builds of vLLM and SGLang may not be at the same commit as the
CUDA-primary builds.  Where version skew exists, it is documented in
Table~\ref{tab:software_versions} in the appendix; performance differences
attributable to feature gaps (rather than hardware) are flagged in the
relevant tables.

\paragraph{No reward-model quality evaluation.}
Blocks C and D (PPO, GRPO, agent RL) use a fixed InternLM2-1.8B reward
model checkpoint.  We measure throughput and the reward signal as seen by
the policy, but do not evaluate downstream model quality (e.g.\ human
preference win rate) on either platform.  The reward model itself is not
trained in this study; its outputs are treated as a black box.

\section{Related Work}
\label{sec:related}

\paragraph{Cross-vendor GPU benchmarks.}
Extensive literature exists on comparing NVIDIA and AMD GPU performance for
deep learning workloads.  \citet{mlperf_training} and \citet{mlperf_inference}~\citep{mlperf2023}
provide standardized training and inference benchmarks across vendors; however,
MLPerf focuses on a fixed set of reference models (ResNet, BERT, GPT-3,
Stable Diffusion) and does not cover the full spectrum of LLM training
algorithms.  Several independent studies have compared
ROCm and CUDA on transformer workloads using frameworks such as PyTorch and
JAX, finding convergent results on most standard operations but noting gaps
in operator coverage and compilation overhead.  Our work extends this by
covering \emph{every stage} of the LLM development pipeline rather than a
single training or inference workload.

\paragraph{Framework-level portability.}
PyTorch's \texttt{torch.accelerator} and the \texttt{torch.amp} unified API
represent the framework's answer to vendor portability~\citep{paszke2019pytorch}.
Earlier work on OpenCL-based portability (e.g.\ PlaidML) did not achieve
the performance parity of native vendor paths.  The ROCm port of PyTorch,
maintained by AMD, has reached a level of maturity where the same Python
source compiles and runs correctly on both CUDA and ROCm; our methodology
exploits this via the thin \texttt{DeviceCtx} wrapper.

\paragraph{Distributed training systems.}
Megatron-LM~\citep{narayanan2021efficient} and DeepSpeed~\citep{rajbhandari2020zero} are the
dominant frameworks for large-scale distributed LLM training, with explicit
AMD/ROCm support added in recent releases.  Both frameworks have published
cross-vendor results, but focus on tensor-parallel and pipeline-parallel
regimes (billions of parameters) that are out of scope for MiniMind's
DDP-only, single-node setup.

\paragraph{RLHF and alignment training.}
The PPO-based RLHF paradigm introduced by \citet{ouyang2022training} and the DPO
simplification of \citet{rafailov2023direct} have been implemented in frameworks such as
TRL and OpenRLHF, some of which have begun adding ROCm support.  MiniMind's
hand-written RLHF trainers (PPO, GRPO, DPO) provide an unmediated comparison
that is not filtered through framework-level abstractions.

\paragraph{MoE model training.}
Mixture-of-experts scaling was studied extensively in~\citet{gshard,switch_transformer}.
Vendor-specific MoE kernels (e.g.\ Megablocks~\citep{megablocks}) have been
developed to reduce token-dispatch overhead on CUDA.  ROCm equivalents via
Composable Kernel are an active area of development.  Our study uses
MiniMind's pure-PyTorch MoE dispatch, which avoids custom kernels and
therefore measures the \emph{baseline} cross-vendor MoE performance
achievable without vendor-specific tuning.

\paragraph{Inference engine benchmarks.}
vLLM~\citep{kwon2023efficient} introduced PagedAttention for efficient KV-cache management;
SGLang~\citep{zheng2024sglang} introduced RadixAttention for prefix sharing.  Both
engines maintain ROCm branches.  Benchmarks comparing vLLM and SGLang on
CUDA exist~\citep{zheng2024sglang}, but systematic cross-vendor comparison using
identical converted weights on the same model is, to our knowledge, not
previously published.

\paragraph{Small models as research vehicles.}
The use of small, reproducible models for algorithm research has precedent in
works such as nanoGPT~\citep{nanogpt} and TinyLlama~\citep{tinyllama}.
MiniMind extends this tradition by including the full post-training and RLHF
pipeline as first-class components, making it suitable for cross-vendor
algorithm coverage studies.

\appendix

\section{Reproducibility Appendix}
\label{sec:appendix}

\subsection{Software Version Matrix}
\label{sec:appendix:versions}

\begin{table}[h]
\centering
\caption{Pinned software versions for all experiments.  ROCm and CUDA wheels
are not interchangeable; the \texttt{torch} version strings differ only in the
\texttt{+cu12x} / \texttt{+rocm7.1} suffix.}
\label{tab:software_versions}
\small
\begin{tabular}{lll}
\toprule
\textbf{Component} & \textbf{NVIDIA (H100)} & \textbf{AMD (MI350)} \\
\midrule
OS                 & Ubuntu 22.04           & Ubuntu 22.04 \\
Kernel             & \PH{}                  & \PH{} \\
Driver             & \PH{} (CUDA 12.x)      & ROCm \PH{} \\
\texttt{torch}     & \texttt{2.10.0+cu12x}  & \texttt{2.10.0+rocm7.1} \\
\texttt{torchvision} & \PH{}                & \PH{} \\
\texttt{transformers} & \PH{}               & \PH{} \\
\texttt{tokenizers} & \PH{}                 & \PH{} \\
\texttt{vllm}      & \PH{} (commit \PH{})   & \PH{} (commit \PH{}) \\
\texttt{sglang}    & \PH{} (commit \PH{})   & \PH{} (commit \PH{}) \\
\texttt{aiohttp}   & \PH{}                  & \PH{} \\
\texttt{swanlab}   & \PH{}                  & \PH{} \\
\texttt{triton}    & \PH{}                  & \PH{} (ROCm Triton) \\
MiniMind commit    & \PH{}                  & same \\
\bottomrule
\end{tabular}
\end{table}

\subsection{Environment Setup}
\label{sec:appendix:setup}

\paragraph{NVIDIA (H100).}

\begin{verbatim}
# Install PyTorch 2.10 CUDA 12.x wheel
pip install torch==2.10.0+cu12x torchvision \
    --index-url https://download.pytorch.org/whl/cu12x

# Install remaining dependencies
pip install transformers tokenizers aiohttp swanlab \
            fastapi uvicorn requests

# Install inference engines (CUDA)
pip install vllm==<PIN>
pip install sglang[all]==<PIN>
\end{verbatim}

\paragraph{AMD (MI350).}

\begin{verbatim}
# ROCm 7.1 must be pre-installed via AMD's repo
# Verify: rocm-smi --version

# Install PyTorch 2.10 ROCm wheel
pip install torch==2.10.0+rocm7.1 torchvision \
    --index-url https://download.pytorch.org/whl/rocm7.1

# Same Python dependencies
pip install transformers tokenizers aiohttp swanlab \
            fastapi uvicorn requests

# Install inference engines (ROCm)
# vLLM ROCm: build from source or use AMD-provided wheel
pip install vllm==<PIN_ROCM>
pip install sglang[all]==<PIN_ROCM>

# Optional: override GFX version for RDNA3 dev machines
# (not needed on MI350 gfx942)
# export HSA_OVERRIDE_GFX_VERSION=11.0.0
\end{verbatim}

\subsection{Dataset Preparation}
\label{sec:appendix:datasets}

All datasets are gitignored and must be downloaded separately.  Checksums
below allow verification that the correct files are in place before running
any experiment block.

\begin{table}[h]
\centering
\caption{Dataset files, sizes, and SHA-256 hashes (first 8 hex digits).}
\label{tab:datasets}
\small
\begin{tabular}{llrr}
\toprule
\textbf{File} & \textbf{Stage} & \textbf{Size} & \textbf{SHA-256 (first 8)} \\
\midrule
\texttt{dataset/pretrain\_t2t\_mini.jsonl} & Pretrain  & \PH{} & \PH{} \\
\texttt{dataset/sft\_t2t\_mini.jsonl}      & SFT       & \PH{} & \PH{} \\
\texttt{dataset/dpo.jsonl}                 & DPO       & \PH{} & \PH{} \\
\texttt{dataset/lora\_medical.jsonl}       & LoRA      & \PH{} & \PH{} \\
\texttt{dataset/rlaif.jsonl}               & PPO/GRPO  & \PH{} & \PH{} \\
\texttt{dataset/agent\_rl.jsonl}           & Agent RL  & \PH{} & \PH{} \\
\texttt{out/internlm2-1\_8b-reward/}       & PPO/GRPO reward & \PH{} & \PH{} \\
\bottomrule
\end{tabular}
\end{table}

\subsection{Canonical Run Commands}
\label{sec:appendix:commands}

All trainer commands are run from the \texttt{trainer/} directory.  All
evaluation and conversion commands are run from the repo root (or the
\texttt{scripts/} directory as noted).

\subsubsection{Block A: Pretrain}

\begin{verbatim}
# Dense, 1 GPU, bf16, no compile
torchrun --nproc_per_node=1 train_pretrain.py \
    --data_path ../dataset/pretrain_t2t_mini.jsonl \
    --hidden_size 768 --num_hidden_layers 8 \
    --epochs 2 --batch_size 32 --dtype bfloat16 \
    --use_compile 0 --use_wandb 0 \
    --save_weight pretrain 2>&1 | tee ../paper/logs/block_a/<VENDOR>_dense_bf16_g1_nocompile.log

# Dense, 8 GPUs, bf16, compile
torchrun --nproc_per_node=8 train_pretrain.py \
    --data_path ../dataset/pretrain_t2t_mini.jsonl \
    --hidden_size 768 --num_hidden_layers 8 \
    --epochs 2 --batch_size 32 --dtype bfloat16 \
    --use_compile 1 --use_wandb 0 \
    --save_weight pretrain 2>&1 | tee ../paper/logs/block_a/<VENDOR>_dense_bf16_g8_compile.log

# MoE, 8 GPUs, bf16
torchrun --nproc_per_node=8 train_pretrain.py \
    --data_path ../dataset/pretrain_t2t_mini.jsonl \
    --hidden_size 768 --num_hidden_layers 8 \
    --use_moe 1 --epochs 2 --batch_size 32 \
    --dtype bfloat16 --use_compile 0 --use_wandb 0 \
    --save_weight pretrain_moe 2>&1 | tee ../paper/logs/block_a/<VENDOR>_moe_bf16_g8.log
\end{verbatim}

\subsubsection{Block B: Post-training}

\begin{verbatim}
# SFT, 8 GPUs
torchrun --nproc_per_node=8 train_full_sft.py \
    --from_weight pretrain \
    --data_path ../dataset/sft_t2t_mini.jsonl \
    --hidden_size 768 --num_hidden_layers 8 \
    --epochs 1 --batch_size 16 --dtype bfloat16 \
    --save_weight full_sft 2>&1 | tee ../paper/logs/block_b/<VENDOR>_sft_g8.log

# LoRA, 1 GPU
torchrun --nproc_per_node=1 train_lora.py \
    --from_weight full_sft \
    --data_path ../dataset/lora_medical.jsonl \
    --hidden_size 768 --lora_name lora_medical \
    --dtype bfloat16 2>&1 | tee ../paper/logs/block_b/<VENDOR>_lora_g1.log

# Distillation: MoE teacher -> dense student
torchrun --nproc_per_node=8 train_distillation.py \
    --teacher_weight pretrain_moe --student_weight full_sft \
    --hidden_size 768 --use_moe 1 --dtype bfloat16 \
    --alpha 0.5 --temperature 1.5 \
    2>&1 | tee ../paper/logs/block_b/<VENDOR>_distill_g8.log
\end{verbatim}

\subsubsection{Block C: Preference Alignment}

\begin{verbatim}
# DPO
torchrun --nproc_per_node=1 train_dpo.py \
    --from_weight full_sft \
    --data_path ../dataset/dpo.jsonl \
    --hidden_size 768 --dtype bfloat16 \
    --beta 0.15 --learning_rate 4e-8 \
    2>&1 | tee ../paper/logs/block_c/<VENDOR>_dpo_g1.log

# GRPO, torch rollout engine, num_generations=4
torchrun --nproc_per_node=8 train_grpo.py \
    --from_weight full_sft \
    --data_path ../dataset/rlaif.jsonl \
    --hidden_size 768 --dtype bfloat16 \
    --num_generations 4 --rollout_engine torch \
    --reward_model_path ../out/internlm2-1_8b-reward \
    2>&1 | tee ../paper/logs/block_c/<VENDOR>_grpo_torch_g4.log

# PPO, sglang rollout engine
torchrun --nproc_per_node=8 train_ppo.py \
    --from_weight full_sft \
    --data_path ../dataset/rlaif.jsonl \
    --hidden_size 768 --dtype bfloat16 \
    --rollout_engine sglang \
    --reward_model_path ../out/internlm2-1_8b-reward \
    2>&1 | tee ../paper/logs/block_c/<VENDOR>_ppo_sglang_g8.log
\end{verbatim}

\subsubsection{Block D: Agent RL}

\begin{verbatim}
torchrun --nproc_per_node=8 train_agent.py \
    --from_weight full_sft \
    --data_path ../dataset/agent_rl.jsonl \
    --hidden_size 768 --dtype bfloat16 \
    --reward_model_path ../out/internlm2-1_8b-reward \
    2>&1 | tee ../paper/logs/block_d/<VENDOR>_agent_g8.log
\end{verbatim}

\subsubsection{Block E: Inference}

\begin{verbatim}
# Convert to Qwen3 HF format (run from repo root)
python scripts/convert_model.py \
    --weight full_sft --hidden_size 768 \
    --out_path out/qwen3_converted/full_sft_blockb_dense_g8

# HF-native single-stream benchmark
python eval_llm.py --weight full_sft --hidden_size 768 \
    --show_speed 1 --max_new_tokens 512

# Concurrent benchmark via OpenAI-compatible server
cd scripts && python serve_openai_api.py \
    --weight full_sft --hidden_size 768 --port 8998 &
sleep 5
python ../paper/experiments/bench_client.py \
    --port 8998 --concurrency 32 --num_requests 200 --max_tokens 512

# vLLM serve (Qwen3-converted weights)
python -m vllm.entrypoints.openai.api_server \
    --model out/qwen3_converted/full_sft_blockb_dense_g8 \
    --port 8999 --tensor-parallel-size 1 &
sleep 10
python paper/experiments/bench_client.py \
    --port 8999 --concurrency 32 --num_requests 200 --max_tokens 512

# Verification: greedy parity (Gate 6)
python paper/experiments/verify_greedy_parity.py \
    --native_weight full_sft_blockb_dense_g8 \
    --conv_path out/qwen3_converted/full_sft_blockb_dense_g8 \
    --server_url http://127.0.0.1:8999 --num_prompts 20

# Tool-call eval
python scripts/eval_toolcall.py \
    --weight full_sft --hidden_size 768
\end{verbatim}

\subsection{Verification Harness}
\label{sec:appendix:verify}

Before reporting any results, run the 6-gate verification harness:

\begin{verbatim}
# For CUDA / H100
bash paper/experiments/verify_all.sh cuda

# For ROCm / MI350
bash paper/experiments/verify_all.sh rocm
\end{verbatim}

Expected terminal output on a clean run:

\begin{verbatim}
[PASS] Gate 1: All banners valid
[PASS] Gate 2: Convergence parity within 2%
[PASS] Gate 3: MoE aux_loss values are sane
[PASS] Gate 4: Checkpoint round-trip passed
[PASS] Gate 5: All pytest tests passed
[PASS] Gate 6: Greedy parity passed
==================================
Verification summary for <vendor>
  PASSED: 6
  FAILED: 0
==================================
All gates passed. Proceed with experiments.
\end{verbatim}

\subsection{Expected Device Banner}
\label{sec:appendix:banner}

Every training run emits a device banner at initialization via
\texttt{print\_device\_info()} (\texttt{trainer/device\_utils.py}).
The banner format is:

\begin{verbatim}
[device] NVIDIA H100 SXM5 80GB | CUDA 12.x | torch 2.10.0+cu12x
         rank=0/8 | bf16=True | dist_backend=nccl
         mem_total=81920 MiB
\end{verbatim}

For AMD:

\begin{verbatim}
[device] AMD Instinct MI350 | ROCm 7.1 | torch 2.10.0+rocm7.1
         rank=0/8 | bf16=True | dist_backend=nccl
         mem_total=<PH> MiB
\end{verbatim}

Gate~1 of the verification harness greps all Block~A logs for the
\texttt{[device]}, \texttt{bf16=True}, and \texttt{dist\_backend=nccl}
strings to confirm that each log corresponds to the intended platform.

\subsection{MFU Calculation Details}
\label{sec:appendix:mfu}

Model FLOPs Utilization is computed as:
\[
  \text{MFU} = \frac{6 \times P \times T_{\text{step}} / t_{\text{step}}}%
                    {F_{\text{peak}}}
\]
where $P$ is the number of active parameters (equal to total parameters for
dense models; equal to $P_{\text{active}}$ for MoE), $T_{\text{step}}$ is
the number of tokens per step (batch size $\times$ sequence length $\times$
number of GPUs), $t_{\text{step}}$ is the wall-clock step time in seconds,
and $F_{\text{peak}}$ is the vendor-advertised peak throughput in FLOPs/s
(H100: 989\,TFLOPs/s bf16; MI350: \PH{}\,TFLOPs/s bf16).

For the dense MiniMind-768 model:
$P_{\text{dense}} \approx 64.7\times10^6$ parameters.

For the MoE MiniMind-768 model with 4 experts, 1 active, and 1 shared:
$P_{\text{active}} = P_{\text{base}} + P_{\text{shared}} + P_{\text{1 routed expert}}
\approx \PH{}\times10^6$ parameters.

The \texttt{paper/experiments/compute\_mfu.py} script automates this
calculation from training logs and \texttt{MiniMindConfig} parameters.

\subsection{Metric Collection}
\label{sec:appendix:metrics}

After all experiment logs are collected:

\begin{verbatim}
# Parse all logs, compute MFU, aggregate metrics
python paper/experiments/collect_metrics.py \
    --log_dir paper/logs \
    --out_dir paper/tables

# Outputs:
#   paper/tables/metrics.json  — full structured results
#   paper/tables/metrics.csv   — flat CSV for spreadsheet import
\end{verbatim}

The JSON schema is:

\begin{verbatim}
{
  "block_a": {
    "<vendor>_<tag>": {
      "tokens_per_s": float,
      "mfu_pct": float,
      "peak_mem_gb": float,
      "epoch_min": float,
      "final_loss": float
    }, ...
  },
  "block_e": {
    "<vendor>_<engine>_<concurrency>": {
      "ttft_p50_ms": float,
      "ttft_p95_ms": float,
      "itl_p50_ms": float,
      "tokens_per_s": float,
      "requests_per_s": float
    }, ...
  }
}
\end{verbatim}


% ---------- references ----------
\bibliographystyle{plainnat}
\bibliography{references}

\end{document}
